\documentclass[../../main.tex]{subfiles}

\begin{document}

\section{依测度收敛}\label{section:ch05-依测度收敛}

我们已考虑了函数序列的一致收敛、逐点收敛、几乎处处逐点收敛, 这里我们增加一种收敛模式, 它既与几乎处处逐点收敛有联系, 又与即将介绍的使得积分号下取极限合理的准则有有用的联系.

\begin{definition}[依测度收敛]\label{definition:ch05-依测度}
令 $\{f_{n}\}$ 为 $E$ 上的可测函数序列而 $f$ 是 $E$ 上的可测函数, 使得 $f$ 和每个 $f_{n}$ 在 $E$ 上 a.e. 有限. 序列 $\{f_{n}\}$ 称为在 $E$ 上\textbf{依测度收敛于} $f$, 若对每个 $\eta>0$,
\begin{align*}
\lim_{n\to\infty}m\{x\in E\mid|f_{n}(x)-f(x)|>\eta\}=0.
\end{align*}
\end{definition}

当我们写在 $E$ 上依测度 $\{f_{n}\}\to f$ 时, 是间接地假设 $f$ 和每个 $f_{n}$ 在 $E$ 上是可测的且 a.e. 有限. 观察到若在 $E$ 上一致地 $\{f_{n}\}\to f$, 且 $f$ 是 $E$ 上的实值可测函数, 则在 $E$ 上依测度地 $\{f_{n}\}\to f$, 这是因为对 $\eta>0$, 当 $n$ 充分大时, 集合 $\{x\in E\mid|f_{n}(x)-f(x)|>\eta\}$ 是空的. 然而, 我们也有以下强得多的结果.

\begin{proposition}\label{proposition:ch05-有限测度a.e.依测度}
假设 $E$ 有有限测度. 令 $\{f_{n}\}$ 为 $E$ 上 a.e. 逐点收敛于 $f$ 的可测函数列且 $f$ 在 $E$ 上 a.e. 有限. 则在 $E$ 上依测度地 $\{f_{n}\}\to f$.
\end{proposition}

\begin{proof}
首先观察到 $f$ 是可测的, 这是由于它几乎处处是可测函数序列的逐点极限. 令 $\eta>0$. 为证明依测度收敛, 我们令 $\varepsilon>0$ 且寻找指标 $N$ 使得
\begin{align}
\text{对所有 }n\geqslant N,\quad m\{x\in E\mid|f_{n}(x)-f(x)|>\eta\}<\varepsilon.\label{align:ch05-依测度估计}
\end{align}
Egoroff 定理告诉我们存在 $E$ 的可测子集 $F$ 满足 $m(E\sim F)<\varepsilon$, 使得在 $F$ 上一致地 $\{f_{n}\}\to f$. 因此存在指标 $N$ 使得在 $F$ 上对所有 $n\geqslant N$, $|f_{n}-f|<\eta$. 因此, 对 $n\geqslant N$, $\{x\in E\mid|f_{n}(x)-f(x)|>\eta\}\subseteq E\sim F$, 因而\eqref{align:ch05-依测度估计}对选取的这个 $N$ 成立.
\end{proof}

若 $E$ 有无穷测度, 以上命题是错的. 以下例子表明该命题的逆也是错的.

\begin{example}\label{example:ch05-滑动帐篷}
考虑 $[0,1]$ 的子区间序列 $\{I_{n}\}_{n=1}^{\infty}$, 其最初的一些项如下:
\begin{align*}
[0,1],[0,1/2],[1/2,1],[0,1/3],[1/3,2/3],[2/3,1],\\
[0,1/4],[1/4,1/2],[1/2,3/4],[3/4,1],\cdots
\end{align*}
对每个指标 $n$, 定义 $f_{n}$ 为 $I_{n}$ 的特征函数在 $[0,1]$ 上的限制. 令 $f$ 为在 $[0,1]$ 上恒等于零的函数. 我们宣称依测度地 $\{f_{n}\}\to f$. 事实上, 观察到 $\lim_{n\to\infty}\ell(I_{n})=0$, 这是由于对每个自然数 $m$,
\begin{align*}
\text{若 }n>1+\cdots+m=\frac{m(m+1)}{2},\quad\text{则 }\ell(I_{n})<1/m.
\end{align*}
因此, 对 $0<\eta<1$, 由于 $\{x\in E\mid|f_{n}(x)-f(x)|>\eta\}\subseteq I_{n}$,
\begin{align*}
0\leqslant\lim_{n\to\infty}m\{x\in E\mid|f_{n}(x)-f(x)|>\eta\}\leqslant\lim_{n\to\infty}\ell(I_{n})=0.
\end{align*}
然而, 显然不存在 $[0,1]$ 中的点 $x$ 使得 $f_{n}(x)$ 收敛到 $f(x)$, 这是由于对 $[0,1]$ 的每个点 $x$, 对无穷多个指标 $n$, $f_{n}(x)=1$, 然而 $f(x)=0$.
\end{example}

\begin{theorem}[Riesz]\label{theorem:ch05-Riesz}
若在 $E$ 上依测度地 $\{f_{n}\}\to f$, 则存在在 $E$ 上 a.e. 逐点收敛于 $f$ 的子序列 $\{f_{n_{k}}\}$.
\end{theorem}

\begin{proof}
根据依测度收敛的定义, 存在严格递增的自然数序列 $\{n_{k}\}$ 使得
\begin{align*}
\text{对所有 }j\geqslant n_{k},\quad m\{x\in E\mid|f_{j}(x)-f(x)|>1/k\}<1/2^{k}.
\end{align*}
对每个指标 $k$, 定义
\begin{align*}
E_{k}=\{x\in E\mid|f_{n_{k}}(x)-f(x)|>1/k\}.
\end{align*}
则 $m(E_{k})<1/2^{k}$, 因此 $\sum_{k=1}^{\infty}m(E_{k})<\infty$. Borel-Cantelli 引理告诉我们, 对几乎所有 $x\in E$, 存在指标 $K(x)$, 使得若 $k\geqslant K(x)$, $x\notin E_{k}$, 即
\begin{align*}
\text{对所有 }k\geqslant K(x),\quad|f_{n_{k}}(x)-f(x)|\leqslant 1/k.
\end{align*}
因此
\begin{align*}
\lim_{k\to\infty}f_{n_{k}}(x)=f(x).
\end{align*}
\end{proof}

\begin{corollary}\label{corollary:ch05-依测度积分收敛}
令 $\{f_{n}\}$ 为 $E$ 上的非负可积函数序列. 则
\begin{align}
\lim_{n\to\infty}\int_{E}f_{n}=0\label{align:ch05-积分收敛零}
\end{align}
当且仅当
\begin{align}
\text{在 }E\text{ 上依测度地 }\{f_{n}\}\to f\equiv 0\text{ 且 }\{f_{n}\}\text{ 在 }E\text{ 上是一致可积与紧的}.\label{align:ch05-依测度一致紧}
\end{align}
\end{corollary}

\begin{proof}
首先假设\eqref{align:ch05-积分收敛零}成立. 推论\cref{corollary:ch05-非负收敛}告诉我们 $\{f_{n}\}$ 在 $E$ 上是一致可积与紧的. 为证明在 $E$ 上依测度 $\{f_{n}\}\to 0$, 令 $\eta>0$. 根据 Chebyshev 不等式, 对每个指标 $n$,
\begin{align*}
m\{x\in E\mid f_{n}>\eta\}\leqslant\frac{1}{\eta}\cdot\int_{E}f_{n}.
\end{align*}
因此,
\begin{align*}
0\leqslant\lim_{n\to\infty}m\{x\in E\mid f_{n}>\eta\}\leqslant\frac{1}{\eta}\cdot\lim_{n\to\infty}\int_{E}f_{n}=0.
\end{align*}
于是在 $E$ 上依测度地 $\{f_{n}\}\to 0$.

为证明逆命题, 我们用反证法. 假设\eqref{align:ch05-依测度一致紧}成立但\eqref{align:ch05-积分收敛零}不成立. 则存在某个 $\varepsilon_{0}>0$ 和子序列 $\{f_{n_{k}}\}$ 使得
\begin{align*}
\text{对所有 }k,\quad\int_{E}f_{n_{k}}\geqslant\varepsilon_{0}.
\end{align*}
然而, 根据定理\cref{theorem:ch05-Riesz}, $\{f_{n_{k}}\}$ 的一个子序列在 $E$ 上几乎处处逐点收敛到 $f\equiv 0$ 且该子序列是一致可积与紧的, 从而根据 Vitali 收敛定理, 我们得到与上述 $\varepsilon_{0}$ 的存在性矛盾的结论. 这完成了证明.
\end{proof}

\begin{exercise}\label{exercise:ch05-06}
令在 $E$ 上依测度地 $\{f_{n}\}\to f$ 而 $g$ 是 $E$ 上的可测函数, 它在 $E$ 上 a.e. 有限. 证明在 $E$ 上依测度地 $\{f_{n}\}\to g$ 当且仅当在 $E$ 上 $f=g$ a.e..
\end{exercise}

\begin{exercise}\label{exercise:ch05-07}
令 $E$ 有有限测度, 在 $E$ 上依测度地 $\{f_{n}\}\to f$ 而 $g$ 是 $E$ 上的一个可测函数, 它在 $E$ 上 a.e. 有限. 证明依测度地 $\{f_{n}\cdot g\}\to f\cdot g$, 且用此证明依测度地 $\{f_{n}^{2}\}\to f^{2}$. 由此推出若依测度地 $\{g_{n}\}\to g$, 则依测度地 $\{f_{n}\cdot g_{n}\}\to f\cdot g$.
\end{exercise}

\begin{exercise}\label{exercise:ch05-08}
证明: 若几乎处处逐点收敛用依测度收敛代替的话, Fatou 引理、单调收敛定理、Lebesgue 控制收敛定理以及 Vitali 收敛定理仍然成立.
\end{exercise}

\begin{exercise}\label{exercise:ch05-09}
证明对无穷测度的集合 $E$, 命题\cref{proposition:ch05-有限测度a.e.依测度}不一定成立.
\end{exercise}

\begin{exercise}\label{exercise:ch05-10}
证明在有限测度集上依测度收敛的序列的线性组合也依测度收敛.
\end{exercise}

\begin{exercise}\label{exercise:ch05-11}
假设 $E$ 有有限测度. 令 $\{f_{n}\}$ 为 $E$ 上的可测函数列而 $f$ 在 $E$ 上是可测的, 使得 $f$ 和每个 $f_{n}$ 在 $E$ 上是 a.e. 有限的. 证明在 $E$ 上依测度地 $\{f_{n}\}\to f$ 当且仅当 $\{f_{n}\}$ 的每个子序列有一个进一步的子序列在 $E$ 上 a.e. 逐点收敛于 $f$.
\end{exercise}

\begin{exercise}\label{exercise:ch05-12}
若实数列 $\{a_{j}\}$ 满足对所有 $j$, $|a_{j+1}-a_{j}|\leqslant 1/2^{j}$, 通过证明序列 $\{a_{j}\}$ 是 Cauchy 的, 证明它收敛到某个实数.
\end{exercise}

\begin{exercise}\label{exercise:ch05-13}
$E$ 上的可测函数序列 $\{f_{n}\}$ 称为是依测度 Cauchy 的, 若给定 $\eta>0$ 和 $\varepsilon>0$, 存在指标 $N$ 使得对所有 $m$, $n\geqslant N$,
\begin{align*}
m\{x\in E\mid|f_{n}(x)-f_{m}(x)|\geqslant\eta\}<\varepsilon.
\end{align*}
证明: 若 $\{f_{n}\}$ 是依测度 Cauchy 的, 则存在 $E$ 上的可测函数 $f$ 使得序列 $\{f_{n}\}$ 依测度收敛到它. (提示: 选取严格递增的自然数序列 $\{n_{j}\}$ 使得对每个指标 $j$, 若 $E_{j}=\{x\in E\mid|f_{n_{j+1}}(x)-f_{n_{j}}(x)|>1/2^{j}\}$, 则 $m(E_{j})<1/2^{j}$. 现在用 Borel-Cantelli 引理和前一个习题.)
\end{exercise}

\begin{exercise}\label{exercise:ch05-14}
假设 $m(E)<\infty$. 对于 $E$ 上的两个可测函数 $g$ 和 $h$, 定义
\begin{align*}
\rho(g,h)=\int_{E}\frac{|g-h|}{1+|g-h|}.
\end{align*}
证明在 $E$ 上依测度地 $\{f_{n}\}\to f$ 当且仅当 $\lim_{n\to\infty}\rho(f_{n},f)=0$.
\end{exercise}

\end{document}